% ==============================================================================
% Fichier : lab1_application.tex
% Description : Laboratoire 1 - Audit d'une Application Comptable (APPLICATION)
% ==============================================================================

\chapter{Laboratoire 1 : Audit d'une Application en Service}\label{chap:lab1_application}

Ce laboratoire met en pratique la méthodologie d'audit sur le cas de l'application comptable de l'entreprise LiZbiZ, nommée \textbf{APPLICATION}. L'objectif est de simuler une mission réelle : revue de contrôle interne, tests de données et rédaction de rapport.

\section{Objectifs du TP}
\begin{itemize}
    \item Créer et peupler l'environnement technique (Base de données).
    \item Concevoir les outils d'audit (Grilles de contrôle).
    \item Réaliser des tests de contrôles et des tests substantifs via SQL.
    \item Rédiger le rapport d'audit final.
\end{itemize}

% ------------------------------------------------------------------------------
\section{Étape 1 : Préparation de l'Environnement (Base de Données)}
% ------------------------------------------------------------------------------

Avant d'auditer, il faut disposer du système. Nous allons créer le schéma de la base de données de l'APPLICATION.

\subsection{Création du Schéma Relationnel}

Exécutez le script SQL suivant sur votre SGBD (PostgreSQL ou MySQL).

\begin{lstlisting}[language=SQL, caption=Script de création de la base APPLICATION]
-- Creation des tables

CREATE TABLE UTILISATEUR (
    IdUser INT PRIMARY KEY,
    Nom VARCHAR(50),
    Prenom VARCHAR(50),
    Service VARCHAR(50),
    Profil VARCHAR(20), -- 'Admin', 'Comptable', 'Validation'
    Actif BOOLEAN
);

CREATE TABLE JOURNAL (
    CodeJournal VARCHAR(5) PRIMARY KEY,
    Libelle VARCHAR(50)
);

CREATE TABLE COMPTE (
    NumCompte VARCHAR(10) PRIMARY KEY,
    Libelle VARCHAR(100),
    Classe VARCHAR(2)
);

CREATE TABLE TIERS (
    IdTiers INT PRIMARY KEY,
    Nom VARCHAR(100),
    Type VARCHAR(20), -- 'Client', 'Fournisseur'
    Email VARCHAR(100)
);

CREATE TABLE ECRITURE (
    IdEcriture INT PRIMARY KEY,
    DateEcriture DATE,
    CodeJournal VARCHAR(5),
    NumCompte VARCHAR(10),
    IdTiers INT,
    Libelle VARCHAR(200),
    Debit DECIMAL(15,2),
    Credit DECIMAL(15,2),
    IdUserCreate INT,
    IdUserValid INT, -- Utilisateur qui valide
    Statut VARCHAR(10), -- 'Brouillon', 'Valide'
    FOREIGN KEY (CodeJournal) REFERENCES JOURNAL(CodeJournal),
    FOREIGN KEY (NumCompte) REFERENCES COMPTE(NumCompte),
    FOREIGN KEY (IdTiers) REFERENCES TIERS(IdTiers)
);
\end{lstlisting}

\subsection{Chargement des Données de Test}

Insérez un jeu de données fictives pour simuler l'activité de LiZbiZ. (Voir annexe du PDF source pour un jeu complet ou générez-en un simplifié).

\begin{lstlisting}[language=SQL, caption=Exemple d'insertion]
-- Insertion des utilisateurs
INSERT INTO UTILISATEUR VALUES (1, 'DUPONT', 'Marie', 'Comptabilite', 'Comptable', TRUE);
INSERT INTO UTILISATEUR VALUES (2, 'MARTIN', 'Paul', 'Direction', 'Validation', TRUE);
INSERT INTO UTILISATEUR VALUES (3, 'HACK', 'Bad', 'IT', 'Admin', TRUE);

-- Insertion d'une ecriture suspecte (Fraude potentielle)
INSERT INTO ECRITURE VALUES (1001, '2023-12-15', 'ACH', '607', 50, 'Achat Fictif', 5000, 0, 1, 1, 'Valide');
\end{lstlisting}

% ------------------------------------------------------------------------------
\section{Étape 2 : Évaluation du Contrôle Interne (COSO)}
% ------------------------------------------------------------------------------

L'objectif est de concevoir la "Checklist" d'audit.

\subsection{Domaine : Environnement de Contrôle}

Concevez un questionnaire d'évaluation pour vérifier la culture sécurité du service comptabilité de LiZbiZ.
\textbf{Questions à formuler :}
\begin{itemize}
    \item Existe-t-il une charte informatique signée par les comptables ?
    \item Les utilisateurs connaissent-ils leur profil d'accès ?
    \item Qui a le profil "Admin" sur l'application ?
\end{itemize}

\subsection{Domaine : Activités de Contrôle (Séparation des Tâches)}

C'est le point critique. Vérifiez la matrice RACI.
\textbf{Tâche :} Vérifier si un même utilisateur peut créer ET valider une écriture (Risque de fraude majeur).

\begin{boxlizbiz}
\textbf{Requête SQL de vérification (Audit Technique) :}
Ecrire une requête pour détecter les utilisateurs qui ont créé une écriture et l'ont eux-mêmes validée.
\begin{lstlisting}[language=SQL]
SELECT IdEcriture, Libelle, IdUserCreate 
FROM ECRITURE 
WHERE IdUserCreate = IdUserValid;
\end{lstlisting}
\textbf{Analyse :} Si le résultat n'est pas vide, c'est une Non-Conformité Majeure (Règle des 4 yeux non respectée).
\end{boxlizbiz}

% ------------------------------------------------------------------------------
\section{Étape 3 : Audit des Données (Tests Substantifs)}
% ------------------------------------------------------------------------------

Nous utilisons le langage SQL comme outil d'audit (CAATs - Computer Assisted Audit Techniques).

\subsection{Test d'Intégrité : Recherche de Doublons}

\textbf{Risque :} Factures payées deux fois.
\textbf{Objectif :} Détecter les écritures strictement identiques (même montant, même date, même tiers, même libellé).

\begin{lstlisting}[language=SQL]
SELECT DateEcriture, IdTiers, Libelle, Sum(Debit) as TotalDebit, Count(*) as NbOcurrences
FROM ECRITURE
GROUP BY DateEcriture, IdTiers, Libelle
HAVING Count(*) > 1;
\end{lstlisting}

\subsection{Test de Cohérence : Equilibre Journal}

\textbf{Risque :} Déséquilibre comptable.
\textbf{Objectif :} Vérifier que Total Débit = Total Crédit pour chaque journal.

\begin{lstlisting}[language=SQL]
SELECT CodeJournal, SUM(Debit) as TotalD, SUM(Credit) as TotalC
FROM ECRITURE
GROUP BY CodeJournal
HAVING SUM(Debit) != SUM(Credit);
\end{lstlisting}

\subsection{Test de Convention : Seuil de Validation}

\textbf{Règle de gestion LiZbiZ :} Toute écriture supérieure à 5000 € doit être validée par le Directeur (Profil 'Validation').
\textbf{Objectif :} Détecter les écritures > 5000 € validées par un simple comptable.

\begin{lstlisting}[language=SQL]
SELECT E.IdEcriture, E.Libelle, E.Debit, U.Nom as Validateur, U.Profil
FROM ECRITURE E
JOIN UTILISATEUR U ON E.IdUserValid = U.IdUser
WHERE E.Debit > 5000 AND U.Profil != 'Validation';
\end{lstlisting}

% ------------------------------------------------------------------------------
\section{Étape 4 : Rapport d'Audit}
% ------------------------------------------------------------------------------

\subsection{Rédaction des Observations}

Sur la base des requêtes exécutées, rédigez une fiche d'observation pour le rapport final (Format : Constat / Risque / Recommandation).

\textbf{Exemple de fiche à compléter :}

\begin{table}[h!]
\centering
\small
\begin{tabular}{|p{3cm}|p{9cm}|}
    \hline
    \textbf{Titre} & Non respect de la séparation des tâches (Création/Validation). \\
    \hline
    \textbf{Constat} & La requête SQL a identifié X écritures validées par leur créateur. \\
    \hline
    \textbf{Risque} & Fraude interne sans contrôle (Détournement possible). \\
    \hline
    \textbf{Recommandation} & Configurer l'APPLICATION pour interdire l'auto-validation (Bloquer l'enregistrement si IdUserCreate = IdUserValid). \\
    \hline
\end{tabular}
\caption{Modèle de fiche d'observation}
\end{table}

% ------------------------------------------------------------------------------
\section{Livrables à rendre}
% ------------------------------------------------------------------------------

\begin{enumerate}
    \item Le script SQL de création de la base (Fichier .sql).
    \item Le fichier Excel ou Word contenant les grilles de contrôle interne.
    \item Les scripts SQL utilisés pour les tests (Audit des données).
    \item Le rapport d'audit final (PDF) contenant les observations et le plan d'action.
\end{enumerate}
